% =============================================================
%  Résumé (Français)
% =============================================================

\chapter*{Résumé}
\addcontentsline{toc}{chapter}{Résumé}

Ce projet de fin d'études porte sur la conception et le développement de \textbf{LoomStay}, une plateforme intelligente de services numériques pour établissements hôteliers, réalisée dans le cadre du Master Professionnel en Génie Logiciel à l'Institut Supérieur d'Informatique de Monastir (ISIMM), en partenariat avec la société \textbf{Loomens}.

L'objectif principal est de digitaliser l'ensemble du parcours client en hôtellerie : depuis le check-in numérique sécurisé par QR code, jusqu'à la gestion des réclamations et le suivi des interventions terrain, en passant par la messagerie bilingue en temps réel et les services d'information hôteliers.

Le système est architecturé selon le patron \textbf{SOA (Architecture Orientée Services)}, composé de quatre modules indépendants communicant via des API REST : une application web React intégrant une PWA installable pour l'espace client chambre et des tableaux de bord pour le personnel, un backend Node.js/Express avec base de données PostgreSQL géré par l'ORM Prisma, un service d'intelligence artificielle FastAPI intégrant la classification automatique des réclamations (régression logistique + TF-IDF) et la traduction multilingue via le modèle pré-entraîné NLLB-200, et une application mobile Flutter pour les agents d'intervention terrain.

Les principales contributions de ce projet incluent un workflow complet de réclamations avec classification automatique par IA, un système de scan QR pour le contrôle d'accès et le suivi horodaté des interventions, une messagerie bilingue client-réception en temps réel avec traduction automatique, et un module de gestion des tâches de ménage de chambre avec suivi par la gouvernante générale.

La méthode agile \textbf{Scrum} a été adoptée, avec une planification en \textbf{3 releases} et \textbf{9 sprints}. Les fonctionnalités d'intelligence artificielle ont été intégrées directement dans les sprints fonctionnels concernés, conformément au principe Scrum de livraison d'incréments complets et potentiellement livrables.

\vspace{0.5cm}
\textbf{Mots-clés :} hôtellerie, services numériques, PWA, check-in numérique, réclamations, classification IA, traduction automatique, NLLB-200, QR code, Flutter, React, Node.js, FastAPI, PostgreSQL, Scrum, SOA.

\newpage
